\section{Motivation}
\label{sec:motivation}

\houdini is designed to provide a systematic and reproducible approach to testing container confinement mechanisms. One of its primary goals is to evaluate whether container security features, such as namespaces, cgroups, seccomp, and Linux capabilities are effectively enforcing isolation. Many container security tools rely on theoretical guarantees or static analysis\cite{clair}, but \houdini focuses on empirical validation by executing controlled security tests that verify the correct enforcement of security mechanisms. By doing so, \houdini ensures that container security mechanisms are correctly configured.

\houdini is not only for testing container isolation but also as a means to optimize container configurations. By systematically simulating various tasks within different containers, \houdini can help determine the minimal set of privileges required for a container to perform its intended function. For example, administrators can start with a container configured with broad privileges and then incrementally reduce permissions, such as Linux capabilities, seccomp filters, or resource limits, while monitoring task success. \houdini’s testing framework will identify the point at which functionality is maintained while unnecessary privileges are removed, thereby pinpointing the least privilege configuration that still allows the container to operate effectively. Minimizing the privileges granted to a container reduces its attack surface, making it more difficult for an attacker to exploit vulnerabilities.

Although runtime monitoring and anomaly detection, such as with eBPF-based tracing\cite{findlay2021bpfcontain} can provide immediate insights into container behavior, our goal isn't to implement an exhaustive, continuous surveillance system. Instead, we aim to focus on auditing and validating the underlying isolation mechanisms to ensure they are correctly configured and consistently enforced.

\subsection{Example Applications of \houdini}
%The following are additional goals of \houdini :
%\todo{i don't think these are goals, but instead appear to be applications.}
A system like \houdini can be used in the following ways: 

\noindent\textbf{Identify gaps between intended confinement policies and actual enforcement:}
There is often a discrepancy between the security policies that administrators intend to enforce and what is actually upheld at runtime. These discrepancies generally fall into two categories: implicit policy problems and conflicting policy configurations.

An implicit policy problem happens when an administrator doesn't explicitly set a security rule, and the system's default behavior is used instead. These defaults might lead to behavior the administrator did not expect. For example, Docker containers are, by default, connected to a shared internal bridge network. While this topology typically prevents internet access, it still allows containers to communicate with each other. An administrator might see that outbound traffic from a container is blocked and mistakenly believe the container is fully isolated from other networks. However, internal communication between containers will still be allowed unless it's explicitly disabled. Thus, this default behavior can create a false sense of security and introduce unexpected risks.

A conflicting policy configuration, on the other hand, occurs when a security policy is correctly specified but is later invalidated by other conflicting configurations. For example, an administrator might launch a container with \texttt{--network=none}, expecting it to be fully isolated from all network communication. However, if the container is also inadvertently given elevated privileges such as \texttt{--cap-add=NET\_ADMIN} or \texttt{--privileged}, it can reconfigure its own network interfaces and regain connectivity. With the \texttt{NET\_ADMIN} capability, processes inside the container gain access to low-level networking operations, such as creating virtual interfaces, assigning IP addresses, or modifying routing tables. A process could use commands like \texttt{ip link set eth0 up} and \texttt{dhclient eth0} to bring up a previously disabled interface and request an IP address from the host’s network. In the case of a \texttt{--privileged} container, even more capabilities are granted, allowing near-complete access to the host's kernel features, including the creation of new network namespaces or interfaces, effectively bypassing the isolation intended by \texttt{--network=none}.

\houdini can help identify these gaps by systematically testing the behavior of the container against the intended security policies, as demonstrated in case study 2 in \cref{sec:evaluation}.

\noindent\textbf{Challenge misleading security assumptions with empirical testing:}
Container security is often discussed in terms of theoretical guarantees, with many security best practices based on assumptions rather than direct testing. For instance, documentation may state that a container running in unprivileged mode cannot escalate to root, but \houdini actively tests whether that is actually the case under different container configurations. By running real-world security validation instead of relying solely on security claims, \houdini helps challenge misleading assumptions and exposes where security mechanisms fail in practice. This is particularly important as container escapes and privilege escalation exploits continue to emerge, demonstrating that security controls do not always work as expected. By grounding security claims in measurable results rather than theory, \houdini helps researchers and practitioners develop more robust container security strategies.

\noindent\textbf{Detect security drift over time:} One of the challenges in maintaining secure containerized environments is configuration drift. Security settings applied during container initialization may be inadvertently changed through image updates, runtime automation, or infrastructure changes. \houdini can be set up to periodically re-evaluate security policies, making it possible to detect when the system diverges from its originally intended security posture. This allows for the timely remediation of regressions that could otherwise go unnoticed. Configuration drift does not only arise from explicit policy changes, but also from underlying updates to system components such as the Linux kernel, Docker, or runc. Even when container configurations and security policies remain unchanged, an upgrade to the kernel or container runtime can subtly alter the enforcement of isolation primitives, security defaults, or available capabilities. For example, a new kernel version may modify the behavior of namespaces, cgroups, or seccomp filters; a Docker update may introduce new features or alter default permissions.

\subsection{Design Goals}
In order to support the use cases mentioned above, we designed \houdini with the following goals:

\begin{enumerate}
  \item \label{dg:repro}\textit{Reproducible Results.} A key goal of \houdini is to ensure that tests yield consistent and repeatable results. By controlling system variables and maintaining a structured testing environment, \houdini allows researchers and practitioners to compare results across different system configurations and security policies.

  \item \label{dg:container}\textit{Separation from the Host.} \houdini is designed to run tests inside a controlled containerized environment within a virtual machine (VM), ensuring that even if a trick triggers a security vulnerability, it does not compromise the host system running the tests. This design choice enhances the safety of security evaluations, preventing unintended side effects on the underlying infrastructure.

  \item \label{dg:test}\textit{Test Case Expressiveness.} \houdini tricks should be maximally expressive, such that some combination of steps can be used to achieve and test any desired result. It should be possible to define a new \houdini trick and modify existing tricks without modifying the \houdini binary. Moreover, it should be clear from reading a defined trick precisely what steps are involved, the consequences of each step passing or failing, and the overall nature of the exploit being tested.

  \item \label{dg:failures}\textit{Focus on Observing Failures, Not Preventing Them.} \houdini is built to identify security weaknesses rather than defend against attacks.
\end{enumerate}
